\documentclass[10pt,a4paper]{article}
\usepackage[T1]{fontenc}
\usepackage{lmodern}
\usepackage{geometry}
\geometry{a4paper,total={170mm,257mm},left=20mm,top=20mm}
\usepackage{amsmath,amssymb,amsthm,mathtools}
\usepackage{microtype,booktabs,array,xcolor}
\usepackage[numbers,sort&compress]{natbib}
\usepackage{xurl}
\usepackage[colorlinks=true,linkcolor=blue!55!black,citecolor=blue!55!black,urlcolor=blue!55!black]{hyperref}
\hypersetup{pdftitle={The Entropy Law for the Asymptotic Subrank of Dicke Tensors},pdfauthor={Wenqi Zhu; Chunlin Liu; Kai Zou; Ping Nie}}
\newtheorem{theorem}{Theorem}[section]
\newtheorem{proposition}[theorem]{Proposition}
\newtheorem{lemma}[theorem]{Lemma}
\newtheorem{corollary}[theorem]{Corollary}
\theoremstyle{definition}
\newtheorem{definition}[theorem]{Definition}
\theoremstyle{remark}
\newtheorem{remark}[theorem]{Remark}
\numberwithin{equation}{section}
\newcommand{\C}{\mathbb C}
\newcommand{\Q}{\mathbb Q}
\newcommand{\F}{\mathbb F}
\newcommand{\E}{\mathbb E}
\newcommand{\Prob}{\mathbb P}
\newcommand{\Qmon}{Q_{\mathrm M}}
\newcommand{\Qt}{\widetilde Q}
\newcommand{\Qmt}{\widetilde Q_{\mathrm M}}
\DeclareMathOperator{\HH}{H}
\DeclareMathOperator{\codim}{codim}
\DeclareMathOperator{\sr}{sr}
\DeclareMathOperator{\rank}{rank}
\DeclareMathOperator{\spanQ}{span_{\Q}}
\title{The Entropy Law for the Asymptotic Subrank of Dicke Tensors}
\author{%
Wenqi Zhu\thanks{Corresponding author.}\\
{\small Mathematical Institute, Woodstock Road, University of Oxford,}\\
{\small Oxford, UK, OX2 6GG}\\
{\small\texttt{wenqi.zhu@maths.ox.ac.uk}}\\[1ex]
Chunlin Liu\\
{\small School of Mathematical Sciences,}\\
{\small University of Science and Technology of China, Hefei, China}\\[1ex]
Kai Zou\\
{\small Mathematical \& Engineering Sciences, King's College London, UK}\\
{\small\texttt{kz@netmind.ai}}\\[1ex]
Ping Nie\\
{\small David R. Cheriton School of Computer Science, University of Waterloo}\\
{\small\texttt{ping.nie@uwaterloo.ca}}}
\date{}

\renewcommand{\baselinestretch}{1.15}
\setlength{\emergencystretch}{2em}
\begin{document}
\maketitle
\begin{abstract}
Let $\mathbf d=(d_1,\ldots,d_n)$ be any tuple of positive integers,
where $n\geq2$, and put $k=\sum_a d_a$ and $p_a=d_a/k$.
The Dicke tensor $D_{\mathbf d}\in(\C^n)^{\otimes k}$ is the sum of
all basis tensors with these multiplicities. We prove that its ordinary
asymptotic subrank and its asymptotic subrank under coordinate restrictions
both equal $\exp\HH(p)$, where $\HH$ is Shannon entropy with natural
logarithms. All $k$ tensor factors remain separate. This proves the
all-multiplicity entropy formula for Dicke tensors. The lower bound
combines rank-dependent hashing and global isolation with an entropy
inequality for finite distributions invariant under an irreducible
linear action. A full joint-type class supplies the candidates, while
every edge of the original support is included in the collision count.
A slice cover gives the finite-power upper bound
$Q(D_{\mathbf d}^{\boxtimes N})\leq k\exp(N\HH(p))$ for every
$N\geq1$, including arbitrary complex-linear restriction maps.
The proof includes the average-free construction, the isolation argument,
and the passage from an unbounded prime-indexed subsequence to the
asymptotic limit. Finite exact checks accompany the written proof.
\end{abstract}

\section{The main theorem}\label{sec:main}

Fix integers $n\geq2$ and $d_1,\ldots,d_n\geq1$, and set
\[
 k=\sum_{a=1}^n d_a,\qquad
 p_a=\frac{d_a}{k},\qquad
 \HH(p)=-\sum_{a=1}^n p_a\log p_a.
\]
All logarithms are natural unless a base is indicated; in particular,
$\HH_2(p)=\HH(p)/\log2$. Write $e_1,\ldots,e_n$ for the standard
basis of $\C^n$. Let $\mathcal W_{\mathbf d}$ be the set of words
$(a_1,\ldots,a_k)\in\{1,\ldots,n\}^k$ in which symbol $a$ occurs
exactly $d_a$ times. The \emph{Dicke tensor} is
\begin{equation}\label{eq:dicke}
 D_{\mathbf d}
   =\sum_{(a_1,\ldots,a_k)\in\mathcal W_{\mathbf d}}
        e_{a_1}\otimes\cdots\otimes e_{a_k}.
\end{equation}
Under the usual symmetric-tensor interpretation it corresponds to a
nonzero scalar multiple of the monomial $x_1^{d_1}\cdots x_n^{d_n}$.
This normalization does not affect the quantities below.

For a nonzero $k$-tensor $T\in V_1\otimes\cdots\otimes V_k$, its
\emph{subrank} $Q(T)$ is the largest integer $q$ for which linear maps
$A_i:V_i\to\C^q$, chosen independently for the $k$ factors, satisfy
\[
 (A_1\otimes\cdots\otimes A_k)T
   =\Delta_q:=\sum_{j=1}^q e_j\otimes\cdots\otimes e_j.
\]
Relative to specified bases of the $V_i$, write $\Qmon(T)$ for the
largest such diagonal obtainable by deleting local basis vectors and
then independently relabelling and rescaling the retained basis vectors.
This is the \emph{coordinate, or monomial, restriction} convention.
In particular, $\Qmon(T)\leq Q(T)$. For the all-one tensor
\eqref{eq:dicke}, it amounts to finding an induced matching in its support.

The grouped power $T^{\boxtimes N}$ is the ordinary $N$th tensor power
with corresponding factors collected together:
\[
 T^{\boxtimes N}\in V_1^{\otimes N}\otimes\cdots\otimes V_k^{\otimes N}.
\]
There are still exactly $k$ factors. Define
\[
 \Qt(T)=\lim_{N\to\infty}Q(T^{\boxtimes N})^{1/N},\qquad
 \Qmt(T)=\lim_{N\to\infty}\Qmon(T^{\boxtimes N})^{1/N}
\]
when these limits exist. Their existence for $D_{\mathbf d}$ is proved
in Lemma~\ref{lem:limits}.

\begin{theorem}\label{thm:main}
For every $n\geq2$ and every tuple $\mathbf d\in\mathbb Z_{>0}^n$,
over $\C$ and with all $k=\sum_a d_a$ original factors separate,
\begin{equation}\label{eq:main}
 \Qmt(D_{\mathbf d})=\Qt(D_{\mathbf d})
       =\exp\HH(\mathbf d/k)=2^{\HH_2(\mathbf d/k)}.
\end{equation}
Moreover, for every integer $N\geq1$,
\begin{equation}\label{eq:finite-upper}
 Q(D_{\mathbf d}^{\boxtimes N})\leq k\exp\bigl(N\HH(\mathbf d/k)\bigr).
\end{equation}
\end{theorem}

The ordinary restriction maps in the theorem need not be symmetric or
monomial. The assertion is pointwise in each fixed multiplicity tuple;
it supplies neither a convergence rate uniform in $\mathbf d$ nor an
efficiency bound for finding the restrictions. It makes no claim that
finite subrank equals the entropy expression, and concerns neither
tensor rank nor border rank. Merging different original modes is not
part of the statement or proof.

For example, when all $d_a$ are equal, the entropy is $\log n$, so
both asymptotic subranks equal $n$. For $\mathbf d=(k-1,1)$ the value is
\[
 \exp\HH\bigl((k-1)/k,1/k\bigr)
   =\frac{k}{(k-1)^{(k-1)/k}}.
\]
These are specializations of the same formula, with the same separation
of the original tensor factors.

\section{Relation to earlier work}
\label{sec:history}

The entropy formula for Dicke tensors was proposed in the study of
asymptotic multipartite entanglement conversion. Vrana and Christandl
\cite[Section~VI, equation~(40)]{VC2015} suggested that the asymptotic
number of copies of a Dicke tensor needed per two-level GHZ tensor is
$1/\HH_2(p)$, where $p=\mathbf d/k$ and $\HH_2$ denotes entropy with
base-two logarithms. Equivalently, the logarithm of its asymptotic subrank
should be $\HH_2(p)$. Oneto and Ventura explicitly restated the ordinary
subrank formula as Conjecture~4.26 in \cite{OV2025}. Their symmetric-tensor
convention for a monomial differs from our Dicke tensor by a nonzero
scalar, which does not affect subrank. Permuting basis symbols also
identifies the partition convention used in some sources with arbitrary
positive multiplicity tuples. Throughout this paper, all $k$ tensor
factors remain separate; tensor powers group only corresponding factors.

A distinction between restriction and degeneration is needed to compare
our conclusion with Question~1.58 of Christandl, Vrana and Zuiddam
\cite{CVZ2019}. Our $\Qmon$ uses coordinate deletions followed by
independent relabellings and nonzero rescalings in the individual factors.
These are monomial restrictions in their Definition~1.6. Their monomial
subexponent $q_{\mathrm M}$, however, is defined using monomial
\emph{degeneration}: see Definition~1.7, equation~(1.9), and
Definition~1.18. Correspondingly, Proposition~1.19 expresses this
subexponent using monomial \emph{border} subrank. Remark~1.12 explicitly
distinguishes the resulting asymptotic rate from the rate defined using
monomial restrictions; the two are not identified there in general.
Writing superscripts to distinguish their notation, our theorem gives
\begin{equation}
\label{eq:cvz-convention-comparison}
\begin{aligned}
 \HH_2(p)
   &= \log_2 \Qmt(D_{\mathbf d}) \\
   &\le q_{\mathrm M}^{\mathrm{CVZ}}(D_{\mathbf d})
    \le q^{\mathrm{CVZ}}(D_{\mathbf d}) \\
   &= \log_2 \Qt(D_{\mathbf d})=\HH_2(p).
\end{aligned}
\end{equation}
The first comparison uses inclusion of monomial restriction in monomial
degeneration, and the second inclusion in ordinary degeneration. The
ordinary restriction--degeneration equivalence of
\cite[Theorem~1.11]{CVZ2019} justifies the last comparison with ordinary
asymptotic subrank. Thus the theorem also answers their Question~1.58,
without assuming any general equivalence between monomial restriction
and monomial degeneration. This external equivalence is needed only for
this comparison of conventions, not for our proof of the ordinary or
coordinate-restriction formula.

We use the higher-order Coppersmith--Winograd construction of
Christandl, Vrana and Zuiddam
\cite[Theorem~1.52 and Section~3.3]{CVZ2019}. The survival probability
of two colliding edges is
controlled by the rational rank of their coordinate differences. We use
this rank-dependent hashing mechanism and give the needed probability
calculation directly. The accompanying average-free sets also have a
standard antecedent: Vrana and Christandl prove the fixed-length
average-free estimate in \cite[Lemma~10 and Appendix~A]{VC2015}, quoted
as Proposition~3.2 in \cite{CVZ2019}. Our proof includes a no-carry digit
construction on a common Euclidean sphere. We make no novelty claim for
the hashing or average-free mechanisms.

Counting a specified class of candidates while testing isolation against
the entire surviving support likewise has a published precedent.
Vrana and Christandl's weighted isolation lemma
\cite[Lemma~2.3, arXiv version~2]{VC2019} permits an arbitrary nonnegative weight on the
full supported set and bounds its sum over globally isolated vertices
of the collision graph. An indicator weight selects a prescribed joint
type for counting purposes while leaving every possible collider in the
graph. Our argument uses this isolation principle with a fixed full
type and derives the required collision estimate from permutation
symmetry. The selection of a full joint type is not presumed to be a
local tensor restriction: the local restriction is obtained only after
isolation in the full support.

There is a related distinction in the entropy optimization of
\cite[Theorem~1.52]{CVZ2019}. Its feasible pair distributions prescribe
only the $2k$ one-coordinate marginals, not both full support-valued
marginal laws. An entropy bound proved for a uniform full first law
therefore cannot simply be substituted into that optimization. We do
not invoke the theorem as a black box or claim to evaluate its displayed
optimization. Instead, we count partners of each fixed candidate by
exact quotient fibers, and bound their entropy by an invariant-subspace
argument. This supplies the collision estimate needed for arbitrary
positive multiplicities, without assuming separate all-binary or
all-distinct cases.

Finally, diagonal slice rank and entropy upper bounds are established
parts of the slice-rank method; see Tao and Sawin
\cite[Example~5 and Proposition~6]{TaoSawin2016}. We include elementary
proofs of the facts used here. A score assigned to each local word
partitions every support entry into a slice and gives
$Q(D_{\mathbf d}^{\boxtimes N})\le k\exp(N\HH(p))$ for every $N\ge1$.
This upper bound applies to arbitrary complex linear maps and completes
the comparison with the coordinate restrictions constructed below.


\section{Powers, limits and the two-factor case}\label{sec:limits}

Each basis vector in a factor of $D_{\mathbf d}^{\boxtimes N}$ is
indexed by a word of length $N$. A support edge is a $k$-tuple of local
words such that each column has multiplicities $\mathbf d$. A set of
edges is an \emph{induced matching} if they have pairwise distinct local
words in every factor and keeping just their local words leaves no
other support edge. A coordinate restriction to such a set, followed
by relabelling, gives a diagonal of that cardinality.

\begin{lemma}\label{lem:limits}
For $R=Q$ and for $R=\Qmon$, the limit
$\lim_{N\to\infty}N^{-1}\log R(D_{\mathbf d}^{\boxtimes N})$
exists and equals
$\sup_{N\geq1}N^{-1}\log R(D_{\mathbf d}^{\boxtimes N})$.
Consequently, a lower bound along any unbounded subsequence also
bounds this limit from below.
\end{lemma}
\begin{proof}
Keeping a single nonzero entry gives $R(D_{\mathbf d}^{\boxtimes N})\geq1$.
Tensoring two restrictions, and identifying a pair of diagonal output
indices with one index, gives
\begin{equation}\label{eq:supermultiplicativity}
 R(D_{\mathbf d}^{\boxtimes(N+L)})
   \geq R(D_{\mathbf d}^{\boxtimes N})R(D_{\mathbf d}^{\boxtimes L}).
\end{equation}
The flattening along one factor has rank at most $n^N$, cannot increase
under restriction, and has rank $q$ on $\Delta_q$. Thus $R\leq n^N$.
Put $a_N=\log R(D_{\mathbf d}^{\boxtimes N})$ and $a_0=0$.
Then $0\leq a_N\leq N\log n$ and $(a_N)$ is superadditive.
Fix $m\geq1$ and write $N=hm+r$ with $0\leq r<m$. Superadditivity
gives $a_N\geq h a_m+a_r\geq h a_m$, whence
\[
 \liminf_{N\to\infty}\frac{a_N}{N}\geq\frac{a_m}{m}.
\]
Take the supremum over $m$. The reverse upper bound is part of the
definition of that supremum, proving the assertion.
\end{proof}

If $k=2$, positivity and $n\geq2$ force $n=2$ and $\mathbf d=(1,1)$.
The tensor $D_{(1,1)}=e_1\otimes e_2+e_2\otimes e_1$ is a permutation
matrix. Its powers have $\Qmon=Q=2^N$, so \eqref{eq:main} holds.
Henceforth we may assume $k\geq3$ in the lower-bound construction.

\section{An entropy inequality from irreducible invariance}
\label{sec:entropy}

We use finite Shannon entropy with natural logarithms.  For a random
variable $X$ with finite support, write
\[
  \HH(X)=-\sum_x \Prob(X=x)\log\Prob(X=x),
\]
where $0\log 0=0$, and define conditional entropy by averaging over the
conditioning variable.  The chain rule is
$\HH(X,Y)=\HH(Y)+\HH(X\mid Y)$.
We will use nonnegativity of conditional mutual information.  To recall
its elementary basis, for probability vectors $p$ and $q$ the relative
entropy
\[
  \operatorname{KL}(p\Vert q)=\sum_{x:p_x>0}p_x\log\frac{p_x}{q_x}
\]
is nonnegative, with value $+\infty$ if some $p_x>0$ has $q_x=0$.
Otherwise, the inequality $-\log u\geq 1-u$ gives
\[
  \operatorname{KL}(p\Vert q)
  \geq \sum_{x:p_x>0}(p_x-q_x)\geq 0.
\]
Conditional mutual information is the average, over the conditioning
variable, of the relative entropy of a conditional joint law with
respect to the product of its marginals.  Thus
\[
  \operatorname{I}(A;B\mid C)
  =\HH(A\mid C)+\HH(B\mid C)-\HH(A,B\mid C)\geq 0.
\]
In particular, if $C$ is a deterministic function of $A$ and also of $B$,
the chain rule yields
\begin{equation}
  \operatorname{I}(A;B)
  =\HH(C)+\operatorname{I}(A;B\mid C)\geq\HH(C).
  \label{eq:common-quotient-information}
\end{equation}

Let $E$ be a finite-dimensional vector space over $\Q$, and let $X$ be
an $E$-valued random variable with finite support.  For a rational linear
subspace $V\leq E$, the notation $X\bmod V$ denotes its coset label,
which is a finite-valued random variable.  Put
\begin{equation}
  F_X(V)=\HH(X\mid X\bmod V)
        =\HH(X)-\HH(X\bmod V).
  \label{eq:quotient-entropy}
\end{equation}
The last equality holds because the coset label is a function of $X$.
In particular, $F_X(0)=0$ and $F_X(E)=\HH(X)$.

\begin{lemma}[Supermodularity]
\label{lem:supermodular}
For rational linear subspaces $U,V\leq E$,
\begin{equation}
  F_X(U)+F_X(V)
  \leq F_X(U+V)+F_X(U\cap V).
  \label{eq:entropy-supermodularity}
\end{equation}
\end{lemma}

\begin{proof}
Set $A=X\bmod U$, $B=X\bmod V$, and $C=X\bmod(U+V)$.
The pair $(A,B)$ and the label $X\bmod(U\cap V)$ determine each other
on the finite support: two support points have the same pair of labels
if and only if their difference belongs to both $U$ and $V$.
Moreover, $C$ is a deterministic function of each of $A$ and $B$.
Equation~\eqref{eq:common-quotient-information} therefore gives
\[
  \HH(X\bmod U)+\HH(X\bmod V)
  -\HH(X\bmod(U\cap V))
  \geq \HH(X\bmod(U+V)).
\]
Substituting \eqref{eq:quotient-entropy} proves the assertion.
\end{proof}

\begin{lemma}[Irreducible invariance]
\label{lem:invariant}
Suppose that $\dim E=t>0$, a group $G$ acts $\Q$-linearly and
irreducibly on $E$, and $gX$ has the same law as $X$ for every $g\in G$.
Then every rational subspace $V\leq E$ satisfies
\begin{equation}
  F_X(V)\leq \frac{\dim V}{t}\,\HH(X).
  \label{eq:invariant-entropy}
\end{equation}
\end{lemma}

\begin{proof}
A finite set has only finitely many partitions.  The partition of the
support into cosets of $V$ determines $F_X(V)$, so this quantity takes
only finitely many values as $V$ varies.  Since the positive dimensions
also belong to a finite set, the maximum
\[
  \lambda=\max_{0\neq V\leq E}\frac{F_X(V)}{\dim V}
\]
exists.  Choose a maximizing subspace $V$ of largest dimension.
Invariance of the law of $X$ gives $F_X(gV)=F_X(V)$, so $gV$ is another
maximizer for every $g\in G$.  Supermodularity and the definition of
$\lambda$ imply
\begin{align*}
  \lambda\bigl(\dim V+\dim gV\bigr)
  &\leq F_X(V+gV)+F_X(V\cap gV)\\
  &\leq\lambda\bigl(\dim(V+gV)+\dim(V\cap gV)\bigr).
\end{align*}
The bound on the intersection also holds when it is zero because
$F_X(0)=0$.  The two outer expressions coincide by the dimension formula.
Thus $V+gV$ attains the maximum slope as well.  The choice of $V$ forces
$\dim(V+gV)=\dim V$, whence $gV=V$.  Since this holds for every $g$,
irreducibility and $V\neq 0$ imply $V=E$.
Consequently $\lambda=F_X(E)/t=\HH(X)/t$, proving the bound for every
nonzero subspace; the zero subspace satisfies it directly.
The argument also applies when $\HH(X)=0$.
\end{proof}

\section{Permutation fibers and counts over the full support}
\label{sec:fibers}

Fix the positive multiplicity tuple $\mathbf d=(d_1,\ldots,d_n)$, with
$n\geq 2$ and $k=\sum_a d_a\geq 3$.  All constants in the remainder of
the proof may depend on this fixed tuple.  Define
\begin{equation}
  p_a=\frac{d_a}{k},\qquad
  M=\frac{k!}{\prod_{a=1}^n d_a!},\qquad
  h=\log M,\qquad
  \alpha_a=ka-\sum_{c=1}^n c\,d_c.
  \label{eq:support-parameters}
\end{equation}
The integers $\alpha_a$ are distinct.  Let $S\subseteq\mathbb Z^k$ be
the set of tuples containing $\alpha_a$ exactly $d_a$ times for every
$a$.  Replacing the symbol $a$ by $\alpha_a$ is only an injective
labelling of the basis symbols of $D_{\mathbf d}$; it is not an operation
on the tensor.  We have $|S|=M$ and
\[
  \sum_{i=1}^k s_i=0\qquad(s\in S).
\]
Let $Z$ be uniform on $S$.  Each coordinate $Z_j$ has distribution $p$
on the integer labels, so
\[
  \HH(Z)=h,\qquad \HH(Z_j)=\HH(p),\qquad
  \HH(p)=-\sum_{a=1}^n p_a\log p_a.
\]
For $1\leq j\leq k$, put
\begin{equation}
  E_j=\left\{v\in\Q^k:
           \sum_{i=1}^k v_i=0,\ v_j=0\right\},
  \qquad \dim E_j=k-2,
  \qquad \beta=\frac{h-\HH(p)}{k-2}\geq 0.
  \label{eq:collision-space}
\end{equation}
Nonnegativity follows from $h-\HH(p)=\HH(Z\mid Z_j)$.

\begin{lemma}[Entropy in a fixed-coordinate fiber]
\label{lem:fiber-entropy}
If $V\leq E_j$ is a rational subspace of dimension $r$, then
\begin{equation}
  F_Z(V)\leq r\beta.
  \label{eq:fiber-entropy-bound}
\end{equation}
\end{lemma}

\begin{proof}
Since every vector of $V$ has coordinate $j$ equal to zero, the quotient
label $Z\bmod V$ determines $Z_j$.  Condition on $Z_j=\alpha_a$.
The other $k-1$ coordinates are uniformly distributed among permutations
of the multiset obtained by removing one copy of $\alpha_a$; their sum
is $-\alpha_a$.  Set the $j$th coordinate to zero and replace every
other coordinate $Z_i$ by $Z_i+\alpha_a/(k-1)$.  This translates the
conditional support into $E_j$.  The translation preserves differences,
entropy, and the conditional entropy modulo $V$.  The translated law
is invariant under all permutations of the $k-1$ coordinates different
from $j$.

The resulting standard action of the symmetric group on $E_j$ is
irreducible over $\Q$.  Indeed, let $W\leq E_j$ be a nonzero invariant
subspace and choose $v\in W\setminus\{0\}$.  Among the $k-1$ moving
coordinates, two entries $v_u,v_w$ must differ, since they sum to zero.
Subtracting the image of $v$ under their transposition yields
\[
  (v_u-v_w)(e_u-e_w)\in W.
\]
The coefficient is nonzero.  Permuting this difference gives all
coordinate differences among the moving coordinates, which span $E_j$.
Thus $W=E_j$.  The same argument includes $k=3$, when $E_j$ is
one-dimensional.

Lemma~\ref{lem:invariant}, applied to each translated conditional law,
now gives
\[
  \HH(Z\mid Z\bmod V,\ Z_j=\alpha_a)
  \leq\frac{r}{k-2}\,\HH(Z\mid Z_j=\alpha_a).
\]
A conditional law supported at one point has entropy zero and obeys the
same inequality.  Average with weights
$\Prob(Z_j=\alpha_a)=p_a$.  Since $Z\bmod V$ determines $Z_j$, the left
side averages to $F_Z(V)$.  The right side averages to
\[
  \frac{r}{k-2}\,\HH(Z\mid Z_j)
  =\frac{r}{k-2}\bigl(h-\HH(p)\bigr)
  =r\beta,
\]
which proves the claim.
\end{proof}

For $s\in S$ and a rational subspace $V\leq E_j$, define
\[
  C_V(s)=(s+V)\cap S.
\]
Conditioned on a coset, the uniform variable $Z$ is uniform on the
corresponding fiber.  Therefore
\begin{equation}
  F_Z(V)=\frac{1}{M}\sum_{s\in S}\log |C_V(s)|.
  \label{eq:fiber-entropy-identity}
\end{equation}

Let $N$ be a positive multiple of $M$.  Identify a support edge of
$D_{\mathbf d}^{\boxtimes N}$ with its sequence of $N$ support columns,
and define
\[
  \mathcal A_N=
  \left\{I=(I(1),\ldots,I(N))\in S^N:
    \#\{t:I(t)=s\}=N/M\ \text{for every }s\in S\right\}.
\]
We call these edges \emph{candidates}.  They form a counting class;
we do not assume that the full joint type $\mathcal A_N$ can itself be
selected by local linear maps.

For a fixed candidate $I\in\mathcal A_N$, each column of a partner
$J\in S^N$ satisfying $J(t)-I(t)\in V$ can be chosen independently
from $C_V(I(t))$.  Hence the count over the \emph{entire original support}
is exactly
\begin{align}
  \#\bigl\{J\in S^N:J(t)-I(t)\in V\ \text{for all }1\leq t\leq N\bigr\}
  &=\prod_{s\in S}|C_V(s)|^{N/M}\notag\\
  &=\exp\bigl(NF_Z(V)\bigr)
   \leq\exp(Nr\beta),\qquad r=\dim V.
  \label{eq:fiber-count}
\end{align}
No joint-type or marginal-type condition is imposed on $J$.
In particular, \eqref{eq:fiber-count} also bounds the smaller class of
partners whose column differences span exactly $V$.

Finally, the size of the candidate class satisfies
\begin{equation}
  \frac{M^N}{(N+1)^M}
  \leq |\mathcal A_N|
  =\frac{N!}{((N/M)!)^M}
  \leq M^N.
  \label{eq:type-count}
\end{equation}
There are at most $(N+1)^M$ count vectors for words in an $M$-letter
alphabet.  The balanced vector maximizes its multinomial coefficient
when $M$ divides $N$: if two counts differ by at least two, transferring
one occurrence from the larger to the smaller strictly decreases the
factorial denominator.  At a maximum, all counts therefore equal
$N/M$.  Since the multinomial classes partition the $M^N$ words, their
largest class has size at least $M^N/(N+1)^M$.  The upper bound follows
from $\mathcal A_N\subseteq S^N$.
For fixed $\mathbf d$, equation~\eqref{eq:type-count} yields
\[
  \log|\mathcal A_N|=Nh-o(N).
\]

\section{Large average-free sets}
\label{sec:average-free-sets}

The hashing argument uses the following form of the no-carry sphere
construction. Generalized average-free sets of this kind occur in
\cite{VC2015,CVZ2019}. We include the construction and its size estimate
to specify exactly the property needed here.

\begin{lemma}[Average-free sets]
\label{lem:average-free}
Fix an integer $L\geq 2$. For every sufficiently large integer $q$ there
is a set
\[
 B\subseteq\left\{0,\ldots,\left\lfloor\frac{q-1}{L}\right\rfloor\right\}
\]
such that
\[
 x_1+\cdots+x_L=Ly,\qquad x_1,\ldots,x_L,y\in B,
 \quad\Longrightarrow\quad x_1=\cdots=x_L=y,
\]
and
\[
 \log|B|=\log q-O_L\!\left(\sqrt{\log q}\right).
\]
The integer $q$ need not be prime in this lemma.
\end{lemma}

\begin{proof}
Put
\[
 \ell=\left\lfloor\sqrt{\log q}\right\rfloor,
 \qquad
 m=\left\lfloor\frac{(q/(2L))^{1/\ell}}{L}\right\rfloor,
 \qquad b=Lm.
\]
For sufficiently large $q$, we have $\ell\geq 1$, $m\geq 2$, and
$b^\ell\leq q/(2L)$. There are $m^\ell$ vectors in
$\{0,\ldots,m-1\}^{\ell}$ and at most $\ell(m-1)^2+1$ possible values
of their squared Euclidean norms. Choose a norm level $\mathcal C$
with
\[
 |\mathcal C|\geq\frac{m^\ell}{\ell(m-1)^2+1},
\]
and encode its vectors in base $b$:
\[
 B=\left\{\sum_{t=0}^{\ell-1}a_t b^t:
                (a_0,\ldots,a_{\ell-1})\in\mathcal C\right\}.
\]
The encodings are distinct and lie between $0$ and $b^\ell-1$.
Since $b^\ell\leq q/(2L)$, they belong to the required interval.

Suppose that $x_1+\cdots+x_L=Ly$ for elements of $B$, and let
$a^{(1)},\ldots,a^{(L)},c\in\mathcal C$ be their digit vectors.
Every digit sum on either side is at most $L(m-1)<b$. Thus neither
side has carries, and uniqueness of the base-$b$ expansion gives
\[
 a^{(1)}+\cdots+a^{(L)}=Lc.
\]
All these vectors have the same squared norm. Consequently
\[
 \frac1L\sum_{i=1}^L\lVert a^{(i)}-c\rVert^2
 =\frac1L\sum_{i=1}^L\lVert a^{(i)}\rVert^2-\lVert c\rVert^2
 =0.
\]
Each $a^{(i)}$ equals $c$, proving the average-free property.

For the size estimate, write $t=\log q$. Then $\ell\asymp\sqrt t$
and
\[
 \log m=\frac{t-\log(2L)}{\ell}-\log L+O_L(m^{-1})
       =\sqrt t+O_L(1).
\]
In particular $m\to\infty$, and
\[
 \ell\log m=t-\log(2L)-\ell\log L+O_L(\ell/m)
             =t-O_L(\sqrt t).
\]
Moreover,
\[
 \log\bigl(\ell(m-1)^2+1\bigr)
 \leq \log(\ell+1)+2\log m=O_L(\sqrt t).
\]
The lower bound for $|\mathcal C|=|B|$ therefore gives
$\log|B|\geq\log q-O_L(\sqrt{\log q})$. The opposite estimate
$|B|\leq q$ completes the proof.
\end{proof}

\section{Hashes and exact survival probabilities}
\label{sec:hashing}

We retain the notation $S\subseteq\mathbb Z^k$, $M$, $h$, $E_j$,
and $\beta$ from Section~\ref{sec:fibers}, with $k\geq 3$ fixed. In
particular, every column in $S$ has coordinate sum zero. We use the
tight-support hashing construction of \cite{CVZ2019}; the uniform rank
bound and the exact probabilities needed here are established below.

\begin{lemma}[Uniform reduction of ranks]
\label{lem:rank-reduction}
Let
\[
 C=\max_{s,s'\in S}\max_{1\leq i\leq k}|s_i-s'_i|,
 \qquad
 B_0>\max\{k,\,k!\max\{1,C\}^{k}\}.
\]
For every prime $q>B_0$, any integer matrix whose columns belong to
$S-S$ has the same rank over $\Q$ and over $\F_q$. The threshold is
independent of the number of columns.
\end{lemma}

\begin{proof}
If the rational rank is zero, every entry is zero and the assertion
is immediate. Otherwise, let that rank be $r$. There is a nonzero
integer $r\times r$ minor, of absolute value at most
\[
 r!C^r\leq k!\max\{1,C\}^{k}<q.
\]
It remains nonzero modulo $q$, so the rank over $\F_q$ is at least
$r$. Every larger minor is zero as an integer and remains zero
modulo $q$, proving the opposite inequality. The argument depends
only on the entries and the number $k$ of rows.
\end{proof}

Fix a prime $q>B_0$ large enough for Lemma~\ref{lem:average-free}
with $L=k-1$, and fix a positive multiple $N$ of $M$. Choose the set
$B$ from that lemma, viewed also as a subset of $\F_q$ by reduction
modulo $q$. Choose
\[
 v=(v_1,\ldots,v_N)\in\F_q^N,
 \qquad u_1,\ldots,u_{k-1}\in\F_q
\]
independently and uniformly. For an integer-labelled local word
$w=(w_1,\ldots,w_N)$, define
\begin{align*}
 b_i(w)&=u_i+\sum_{t=1}^Nw_t v_t &&(1\leq i<k),\\
 b_k(w)&=\frac{\sum_{i<k}u_i-\sum_{t=1}^Nw_t v_t}{k-1}.
\end{align*}
These expressions are in $\F_q$; the denominator is invertible
because $q>k$.

An original support edge $I\in S^N$ may be described by its columns
$I(t)\in S$ or by its $k$ local words $I_1,\ldots,I_k$. The
zero-sum condition on every column gives
\begin{equation}
 \sum_{i<k}b_i(I_i)=(k-1)b_k(I_k)\qquad\text{in }\F_q.
 \label{eq:tight-hash-relation}
\end{equation}
Keep in each mode exactly the local basis words whose hash belongs
to $B$. This is a coordinate restriction of all $k$ factors. If
$I$ survives, represent each of its hashes by the corresponding
integer in $B$. Both sides of \eqref{eq:tight-hash-relation} then
belong to $[0,q-1]$, so their congruence is an equality of integers.
Lemma~\ref{lem:average-free} implies that all $k$ hashes equal a
single element of $B$. Conversely, this condition implies survival.

For a fixed edge $I$ and a prescribed common hash $b\in B$, each
choice of $v$ permits exactly one offset tuple: set
$u_i=b-\sum_t I_i(t)v_t$ for $i<k$. The last hash then also equals
$b$ by \eqref{eq:tight-hash-relation}. The events for distinct $b$
are disjoint. Of the $q^{N+k-1}$ choices of all hashing variables,
exactly $|B|q^N$ make $I$ survive. Thus
\begin{equation}
 \Prob[I\text{ survives}]=|B|q^{-(k-1)}.
 \label{eq:single-survival}
\end{equation}

Now let distinct original edges $I,J\in S^N$ share a local word,
say $I_j=J_j$. Put
\[
 V=\operatorname{span}_{\Q}
       \{I(t)-J(t):1\leq t\leq N\},
 \qquad r=\dim V.
\]
Every column difference has coordinate sum zero and coordinate $j$
equal to zero. Hence $V\subseteq E_j$ and $r\leq k-2$. Also
$r\geq 1$, since the integer labelling is injective and $I\ne J$.

If both edges survive, their common hash values agree: the shared
local word has the same hash for both. For a prescribed common
$b\in B$, the offsets are uniquely determined by $I$ once $v$ is
chosen. The additional conditions on $J$ are precisely
\[
 \sum_{t=1}^N\bigl(I_i(t)-J_i(t)\bigr)v_t=0,
 \qquad 1\leq i<k.
\]
The rows of the full $k\times N$ difference matrix sum to zero.
Its last row is therefore redundant, over both $\Q$ and $\F_q$.
Deleting that row leaves rank $r$ by
Lemma~\ref{lem:rank-reduction}. The displayed system has exactly
$q^{N-r}$ solutions for $v$, each with one associated offset tuple.
Consequently
\begin{equation}
 \Prob[I\text{ and }J\text{ survive}]
       =|B|q^{-(k-1+r)}.
 \label{eq:pair-survival}
\end{equation}
This argument also applies when the shared mode is $j=k$: its
shared hash still forces equality of the two common values, and
the last row remains redundant. Neither edge is required to have
uniform joint type for \eqref{eq:single-survival} or
\eqref{eq:pair-survival}.

\section{Isolation against the full original support}
\label{sec:isolation}

We now distinguish the counting class $\mathcal A_N$ from the
support $S^N$ of the full tensor power. Let $X$ count surviving
candidates $I\in\mathcal A_N$. Let $Y$ count ordered pairs $(I,J)$
such that $I$ is a surviving candidate, $J\ne I$ is any surviving
original edge, and $I$ and $J$ share a local word in at least one
mode. In particular, no type condition is imposed on $J$.
Equation~\eqref{eq:single-survival} gives
\begin{equation}
 \E X=|\mathcal A_N|\,|B|q^{-(k-1)}.
 \label{eq:candidate-expectation}
\end{equation}

Let $\mathcal V$ be the collection of distinct nonzero rational
spans of subsets of $S-S$ that are contained in at least one $E_j$,
and put $K=|\mathcal V|$. This collection is finite because $S-S$
is finite; it depends on $\mathbf d$, but not on $N$. Every
$V\in\mathcal V$ has dimension between $1$ and $k-2$, and each
pair counted by $Y$ has exactly one such difference span.

Fix $I\in\mathcal A_N$ and $V\in\mathcal V$ of dimension $r$.
The number of partners $J\in S^N$ whose column differences with
$I$ span exactly $V$ is at most the number whose differences all
belong to $V$. By \eqref{eq:fiber-count} and
Lemma~\ref{lem:fiber-entropy}, this latter number is at most
$\exp(Nr\beta)$. We use a mode $j$ with $V\subseteq E_j$ only to
apply that entropy estimate; we do not sum over choices of $j$.
Summing \eqref{eq:pair-survival} over the candidates and their
possible difference spans therefore yields
\begin{equation}
 \E Y\leq \E X
       \sum_{V\in\mathcal V}
          \exp\bigl(N\dim(V)\beta\bigr)q^{-\dim(V)}.
 \label{eq:collision-expectation}
\end{equation}
The larger fiber count is an upper bound even though it also
allows partners of smaller difference span. Each actual pair is
assigned to its unique exact span in this sum, so there is no
extra multiplicity for the shared mode.

For a realization of the hashing variables, call a surviving
candidate \emph{isolated} if it shares no local word with any other
surviving original edge. Denote the set of these candidates by
$\mathcal I$. Every nonisolated candidate contributes at least one
ordered pair to $Y$, and hence
\[
 |\mathcal I|\geq X-Y.
\]
This is the candidate-indicator instance of the weighted
global-isolation principle from \cite{VC2019}. The point needed
for a coordinate restriction is that the comparison uses the
entire original support, as follows.

If $\mathcal I$ is empty, then $X-Y\leq 0$ and the corresponding
lower estimate follows already from $\Qmon\geq 1$. Otherwise,
retain exactly the local words belonging to the edges in
$\mathcal I$. Distinct edges of $\mathcal I$ have distinct local
words in every mode, since they are isolated. Any original edge
remaining after this final coordinate restriction would have
survived the earlier hash restriction. It also shares a retained
word with some $I\in\mathcal I$. If it differed from $I$, it would
contradict the definition of isolation. Thus the remaining support
is exactly $\mathcal I$, and independent relabelling in the $k$
modes gives the diagonal tensor of size $|\mathcal I|$.

The hashing probability space is finite, so some realization has
$X-Y\geq\E(X-Y)$. We have proved
\begin{equation}
 \Qmon\bigl(D_{\mathbf d}^{\boxtimes N}\bigr)\geq\E(X-Y).
 \label{eq:expected-diagonal}
\end{equation}
No local restriction selecting $\mathcal A_N$ was used. Isolation
only within the candidate class would not justify the conclusion:
an edge of $S^N\setminus\mathcal A_N$ could remain on the retained
local words.

\section{The entropy lower bound}
\label{sec:lower-bound}

Fix $\delta>0$ and set $\lambda=\beta+\delta>0$. Choose primes
$q$ tending to infinity, and for sufficiently large $q$ define
\[
 N=N(q)=M\left\lfloor\frac{\log q}{M\lambda}\right\rfloor.
\]
Then $N$ is a positive multiple of $M$, tends to infinity with
$q$, and satisfies
\begin{equation}
 N\lambda\leq\log q<(N+M)\lambda.
 \label{eq:prime-first-subsequence}
\end{equation}
The primes are chosen first and the tensor powers second. Thus
unboundedness of the primes suffices; no estimate for prime gaps
is needed. For all sufficiently large $q$, both the rank threshold
and Lemma~\ref{lem:average-free} apply.

For $r=\dim V\geq 1$, equation~\eqref{eq:prime-first-subsequence}
gives
\[
 \exp(Nr\beta)q^{-r}\leq\exp(-Nr\delta)
                       \leq\exp(-N\delta).
\]
It follows from \eqref{eq:collision-expectation} that
\[
 \E Y\leq\E X\,K\exp(-N\delta).
\]
Combining this inequality with \eqref{eq:candidate-expectation}
and \eqref{eq:expected-diagonal} yields
\begin{equation}
 \Qmon\bigl(D_{\mathbf d}^{\boxtimes N}\bigr)
 \geq\bigl(1-Ke^{-N\delta}\bigr)
        |\mathcal A_N|\,|B|q^{-(k-1)}
 \label{eq:finite-monomial-bound}
\end{equation}
for these powers and all sufficiently large $q$. The prefactor
is positive and tends to one.

All of $\mathbf d,k,M,K,\beta,\delta$ are fixed before $q$ tends
to infinity. By \eqref{eq:type-count},
$\log|\mathcal A_N|=Nh-o(N)$. Lemma~\ref{lem:average-free} gives
$\log|B|=\log q-O_k(\sqrt{\log q})$, and
\eqref{eq:prime-first-subsequence} gives
$\log q=N(\beta+\delta)+O_{\mathbf d,\delta}(1)$.
In particular $\log q=O_{\mathbf d,\delta}(N)$, so the error in
the average-free estimate is $o(N)$. The logarithm of the right
side of \eqref{eq:finite-monomial-bound} is consequently
\begin{align*}
 Nh-(k-2)\log q-o(N)
 &=N\bigl[h-(k-2)(\beta+\delta)\bigr]-o(N)\\
 &=N\bigl[\HH(p)-(k-2)\delta\bigr]-o(N).
\end{align*}
Here we used $(k-2)\beta=h-\HH(p)$.

Lemma~\ref{lem:limits} transfers this lower bound along the
unbounded sequence of powers to the asymptotic limit, giving
\[
 \log\Qmt(D_{\mathbf d})\geq\HH(p)-(k-2)\delta.
\]
Since $\delta>0$ was arbitrary, and every coordinate restriction
is a linear restriction, we conclude that
\begin{equation}
 \Qt(D_{\mathbf d})\geq\Qmt(D_{\mathbf d})
                    \geq\exp\bigl(\HH(p)\bigr).
 \label{eq:monomial-lower}
\end{equation}

\section{The upper bound for arbitrary linear restrictions}\label{sec:upper}

A \emph{slice in mode $i$} is a tensor of the form $a\otimes B$,
with $a$ in the $i$th factor and $B$ in the tensor product of the other
factors, placed in their original order. The \emph{slice rank} $\sr(T)$
is the least number of slices summing to $T$; their distinguished modes
may differ. Applying linear maps independently to the factors takes
a slice to a slice or zero. We first prove the diagonal value needed
to bound ordinary subrank. This fact and the entropy upper-bound method
have precedents in \citet{TaoSawin2016}; the required arguments follow.

\begin{lemma}\label{lem:product-dimension}
For subspaces $U,V\subseteq\C^q$, let $U*V$ be the linear span of
their coordinatewise products. Then
\begin{equation}\label{eq:product-dimension}
 \dim(U*V)\geq\dim U+\dim V-q.
\end{equation}
Consequently, for any $\ell\geq1$ and $U_1,\ldots,U_\ell\subseteq\C^q$,
\begin{equation}\label{eq:iterated-product}
 \dim\operatorname{span}\{u_1*\cdots*u_\ell:u_i\in U_i\}
    \geq q-\sum_{i=1}^{\ell}\codim U_i.
\end{equation}
\end{lemma}
\begin{proof}
Let $A$ be the union of the coordinate supports of vectors in $U$.
There is $u\in U$ nonzero in every coordinate of $A$: each such
coordinate is a nonzero linear function on $U$, and a finite union of
their proper kernels cannot cover a vector space over the infinite
field $\C$. Equivalently, the product of these linear functions is
a nonzero polynomial and is nonzero somewhere. If $U=0$, take $A$
empty and $u=0$.

Coordinatewise multiplication by $u$ on $\C^q$ has kernel of dimension
$q-|A|$. Since $\dim U\leq|A|$, its restriction to $V$ gives
\[
 \dim(U*V)\geq\dim(u*V)
     \geq\dim V-(q-|A|)\geq\dim V+\dim U-q.
\]
Iterate this inequality, using bilinearity of coordinatewise
multiplication to identify the successive spans. The case $\ell=1$
of \eqref{eq:iterated-product} is an equality.
\end{proof}

\begin{lemma}\label{lem:diagonal-slice}
For every $k\geq2$ and $q\geq1$, $\sr(\Delta_q)=q$.
In particular, $Q(T)\leq\sr(T)$ for every $k$-tensor $T$ over $\C$.
\end{lemma}
\begin{proof}
Consider any slice decomposition of $\Delta_q$. For each $i$, let
$r_i$ be the number of its slices in mode $i$, and let $U_i\subseteq\C^q$
annihilate their distinguished coefficient vectors under the bilinear
pairing $\sum_j x_jy_j$. Then $\codim U_i\leq r_i$. Evaluating the
decomposition on $U_1\times\cdots\times U_k$ kills every slice, so
\[
 \sum_{j=1}^q u_{1,j}\cdots u_{k,j}=0
       \qquad (u_i\in U_i).
\]
Thus the span of $U_1*\cdots*U_{k-1}$ lies in $U_k^\perp$, whose
dimension is $\codim U_k$. By \eqref{eq:iterated-product},
\[
 q-\sum_{i<k}\codim U_i\leq\codim U_k,
       \qquad q\leq\sum_i\codim U_i\leq\sum_i r_i.
\]
The defining expression for $\Delta_q$ is a sum of $q$ slices, proving
equality. Monotonicity of slice rank under independent linear maps
now gives $q=\sr(\Delta_q)\leq\sr(T)$ for every diagonal restriction
of $T$, and hence $Q(T)\leq\sr(T)$.
\end{proof}

\begin{proposition}\label{prop:upper}
For all $N\geq1$,
\[
 Q(D_{\mathbf d}^{\boxtimes N})\leq
 \sr(D_{\mathbf d}^{\boxtimes N})\leq k\exp(N\HH(p)).
\]
\end{proposition}
\begin{proof}
For a local word $w=(w_1,\ldots,w_N)\in\{1,\ldots,n\}^N$, set
\[
 s(w)=-\sum_{t=1}^N\log p_{w_t},\qquad
 P(w)=\prod_{t=1}^N p_{w_t}=e^{-s(w)},\qquad
 L_N=\{w:s(w)\leq N\HH(p)\}.
\]
All $p_a$ are positive. Since $\sum_wP(w)=1$ and every $w\in L_N$
has $P(w)\geq e^{-N\HH(p)}$, we have
\begin{equation}\label{eq:score-count}
 |L_N|\leq e^{N\HH(p)}.
\end{equation}
For every support edge $(w^{(1)},\ldots,w^{(k)})$ of the full tensor
power, each column has multiplicities $\mathbf d$. Therefore
\[
 \sum_{i=1}^k s(w^{(i)})
    =-N\sum_{a=1}^n d_a\log p_a=Nk\HH(p).
\]
At least one local word belongs to $L_N$. Assign the edge to the least
such mode. For each $i$ and $w\in L_N$, collect all assigned terms
whose $i$th word is $w$ into one slice with distinguished factor $e_w$.
This partitions the entire support and yields at most $k|L_N|$ slices.
Apply Lemma~\ref{lem:diagonal-slice} and \eqref{eq:score-count}.
\end{proof}

Taking $N$th roots in Proposition~\ref{prop:upper} and using
Lemma~\ref{lem:limits} gives $\Qt(D_{\mathbf d})\leq e^{\HH(p)}$.
Together with the coordinate-restriction lower bound
\eqref{eq:monomial-lower} and $\Qmt\leq\Qt$, this proves
Theorem~\ref{thm:main} for $k\geq3$. The two-factor case was handled
in Section~\ref{sec:limits}. Every restriction in the argument retains
the same $k$ factors.

\section{Exact finite verification}\label{sec:verification}

The proof above establishes the theorem for every multiplicity tuple.
Two accompanying standard-library Python programs provide finite exact
checks of its entropy and hashing interfaces. They use integer and
rational arithmetic, explicit failure exits and negative controls;
no floating-point comparison or numerical tolerance enters the checks.
The finite computations are supporting regressions, not premises of
the universal theorem or an extrapolation to larger multiplicities.

\subsection{Collision subspaces and the entropy inequality}

The program \texttt{check\_entropy.py} enumerates all nonzero
difference-generated collision subspaces for one fixed shared mode
in the five supports shown in Table~\ref{tab:entropy-checks}. The
other modes are equivalent by permutation symmetry. It forms the
quotient fibers and independently checks their cardinalities by exact
difference membership. The zero subspace has the immediate zero-entropy
bound and is omitted.

For a checked support of size $M$, let
$A=\prod_{s\in S}|C_V(s)|$ and $r=\dim V$. Rather than approximate
logarithms in Lemma~\ref{lem:fiber-entropy}, the program checks the
equivalent integer inequality
\begin{equation}\label{eq:integer-entropy-check}
 A^{k(k-2)}k^{kMr}
    \leq M^{kMr}\left(\prod_{a=1}^n d_a^{d_a}\right)^{Mr}.
\end{equation}
Indeed, $\HH(p)=\log k-k^{-1}\sum_a d_a\log d_a$, and multiplying
$M^{-1}\log A\leq r(\log M-\HH(p))/(k-2)$ by $Mk(k-2)$ and
exponentiating gives \eqref{eq:integer-entropy-check}.

\begin{table}[htbp]
\centering
\begin{tabular}{lrr}
\toprule
Multiplicities $\mathbf d$ & Support size $M$ & Nonzero collision subspaces\\
\midrule
$(1,1,1)$ & 6 & 1\\
$(2,1,1)$ & 12 & 7\\
$(1,1,1,1)$ & 24 & 13\\
$(2,2,1)$ & 30 & 191\\
$(3,1,1)$ & 20 & 176\\
\bottomrule
\end{tabular}
\caption{Exact finite checks of the conditional entropy bound.}
\label{tab:entropy-checks}
\end{table}

Each case includes a negative control that decreases the full collision
space's claimed dimension by one while leaving its fibers fixed. The
same inequality validator must reject the resulting false claim.
All five controls reject. This tests selected dimensions and supports;
the invariant-entropy argument, rather than this enumeration, proves
the all-subspace statement.

\subsection{Hashing and isolation in the original support}

The program \texttt{check\_hashing.py} checks 405 rank-one collision-subspace
cases, counting the shared mode as part of the case, across multiplicities
\[
 (2,1),\ (1,1,1),\ (2,2),\ (2,1,1),\ (1,1,1,1),\ (3,1,1),\ (2,2,1).
\]
It also exhausts a complete small hashing experiment. Take
\[
 S=\{(-1,-1,2),(-1,2,-1),(2,-1,-1)\},\quad
 N=3,\quad q=11,\quad B=\{0,1,3\}.
\]
There are $11^5=161{,}051$ assignments of the common hash vector and
two offsets, 27 original support edges, six candidates of uniform joint
type, and 54 ordered candidate/collider pairs sharing a local word.
For every assignment, the program checks equal hashes on surviving
edges, the $X-Y$ isolation inequality, and that retaining the isolated
candidates' vertices leaves exactly those edges in the \emph{full}
original support.

Each edge survives in 3,993 assignments and each listed pair in 363,
matching \eqref{eq:single-survival} and \eqref{eq:pair-survival} exactly.
The totals are $\sum X=23{,}958$ and $\sum Y=19{,}602$, so the
average of $X-Y$ is positive. The largest isolated matching obtained
has two edges. These small numbers test the mechanism, not its
asymptotic rate.

Three further controls reject: a two-point distribution on a line in
dimension two refutes omission of invariance from the entropy-slope
bound; replacing $B$ by $\{0,1,2\}$ violates average-freeness through
$0+2=2\cdot1$; and keeping only candidate edges $(0,0)$ and $(1,1)$
while the original support also contains $(0,1)$ exposes a hidden mixed
edge. The first is a counterexample to dropping an assumption, not an
algorithm for detecting all invariant distributions.

Both checkers pass with byte-identical mathematical output under normal
Python and optimized Python (\texttt{-O}); all eight negative controls
reject in both modes. Reproduction instructions and exact input/output
digests accompany the manuscript. These computations make no claim about
a uniform convergence rate, efficient extraction of asymptotic restrictions,
or formal proof-assistant certification.



\section*{Statements and Declarations}
\paragraph{Funding.}
The research of Wenqi Zhu was supported by the Engineering and Physical
Sciences Research Council under grant EP/Y028872/1, \emph{Mathematical
Foundations of Intelligence: An ``Erlangen Programme'' for AI}.

\paragraph{Competing interests.}
The authors have no competing interests to declare that are relevant to
the content of this article.

\paragraph{Data availability.}
The exact verification scripts and Lean~4 files accompanying
this paper are available at
\url{https://github.com/erenup/dicke-tensor-subrank}.

\paragraph{Use of generative AI and AI-assisted technologies.}
During the preparation of this work, the authors used AI tools (Codex and Claude Code) to generate candidate constructions and code, perform computations, draft mathematical arguments, and improve the readability of the text. After using these
tools, the authors reviewed and edited the content as needed and take full
responsibility for the content of the publication. Finite exact checks (Section~\ref{sec:verification}) and a Lean~4 formalization of ten finite algebraic and combinatorial ingredients of the proof, with its scope stated in the repository above, accompany the written proof.

\bibliographystyle{plainnat}
\bibliography{refs}
\end{document}
